% !TeX root = ../../main.tex
\subsection{Nominal-versus-realized null calibration}
\label{sec:supp-calibration}

In the prespecified continuous-dose simulation, at a nominal 0.05 threshold,
unmoderated BARCS had gene-level type-I error 0.081 and empirical FDP 0.130;
moderation reduced FDP to 0.091 but did not repair the marginal type-I error
(Table~\ref{tab:calibration-diagnostic}).  Official MAGeCK-MLE was slightly
conservative in this realization.  Both the plug-in beta-binomial statistic
and its moderated form therefore require external control diagnostics to
assess calibration.

\begin{table}[htbp]
\centering
\caption{Nominal-versus-realized error in the prespecified continuous-dose
simulation.  Type-I error is the fraction of 160 null genes with $p<0.05$;
empirical FDP is measured among calls at Benjamini--Hochberg FDR 0.05.  This single deterministic simulation provides a diagnostic and cannot
establish a uniform operating characteristic.}
\label{tab:calibration-diagnostic}
\small
\begin{tabular}{lrrrr}
\toprule
Method & Null type I & Power & Calls & Empirical FDP\\
\midrule
Naive binomial $z$ & 0.844 & 1.000 & 174 & 0.770\\
BARCS unmoderated $t$ & 0.081 & 1.000 & 46 & 0.130\\
BARCS moderated $t$ & 0.088 & 1.000 & 44 & 0.091\\
Official MAGeCK-MLE Wald & 0.038 & 0.975 & 41 & 0.049\\
\bottomrule
\end{tabular}
\end{table}

We then simulated all-null beta-binomial screens with $m\in\{4,6,8,12\}$
balanced independent libraries, 100 genes, three or five guides per gene, and
a library total of 100,000.  Mean guide counts were 10, 100, or 1,000.
Dispersion was parameterized by $(N-1)\rho\in\{0,5,60,180\}$, corresponding to
$\rho\in\{0,0.00005,0.00060,0.00180\}$ and spanning the 60--180 range in the
continuous-dose simulation.  Cells combining mean count 10 with
$(N-1)\rho\geq60$ were omitted because fewer than 50 guide fits were
estimable.  Each remaining combination used five prespecified seeds.

The tested group coefficient was zero for every guide.  Guide probabilities
were combined by the directional Stouffer rule used in the benchmarks.  To
test calibration after aggregation, odd-numbered null genes estimated a
signed-normal 95th-tail scale and even-numbered null genes evaluated it. A
focused arm at $m=8$, mean count 100, and $(N-1)\rho\in\{60,180\}$ generated
beta probabilities with a Gaussian copula having within-gene guide correlation
0.4.  Table~\ref{tab:calibration-grid} reports ranges across the applicable
dispersion, abundance, and sample-size cells.  Complete parameter-by-parameter
and seed-level values accompany the manuscript.

Under independent guides, expanding dispersion through the continuous-dose
range did not reproduce the 0.081 gene-level error: the largest five-seed mean
was 0.074, and error did not increase monotonically with dispersion.
Differences in design and averaging across seeds limit comparison with the
continuous-dose result.  By contrast, within-gene guide dependence left
guide-level error near or below 0.05 but raised gene-level error to
0.108--0.262.  Aggregation-matched split-control scaling reduced this range to
0.056--0.112, but did not fully repair the five-guide correlated cells.
Calibration must therefore account for the gene combiner and within-gene
dependence, with the control scale estimated at the same aggregation level as
the reported statistic.

The fraction of converged guide fits with $\widehat\rho=0$ was 11/86,840
(0.013\%) in HT-29, 0/5,999 in IL2RA, and 1,738/280,541 (0.620\%) across the
four Liang cell lines.  Boundary truncation is therefore common in the
deliberately near-binomial grid but rare in the three real-data fits and
appears unlikely to explain their calibration behavior in general.  Limited
independent-library degrees of freedom and the plug-in treatment of dispersion
remain relevant.
